This technical note documents household fuel and co-emissions in the shared PakEmis national inventory, whose common annual core spans July 2015--June 2025. Official fuel deliveries, household surveys and agricultural material flows support separate reference S and alternative L configurations. Reference S uses survey-reconstructed household dung; L assumes all agricultural-ledger diverted dung is burned in households. Conservation establishes diverted material but does not determine its destination, so the unresolved quantity remains explicit.

The reference household module contributes \SI{\SectorRefGHG}{\mega\tonne} \COe{} on AR5 GWP-100 and \SI{\SectorRefPMfine}{\kilo\tonne} primary PM$_{2.5}$ in FY2023--24. FY2024--25 quantities are \SI{\SectorLatestGHG}{\mega\tonne} and \SI{\SectorLatestPMfine}{\kilo\tonne}, respectively. Pollutants are native masses; BC and OC are particulate components. Biomass factors use dry fuel including ash and fossil factors retain their applicable energy and particulate definitions.

Historical survey interpolation, stove/fuel factor transfer and structural dung fate remain important limitations. Conditional uncertainty shares activity across gases and pollutants, with ledger dung following the same sampled herd as manure accounting. Spatial products allocate the supported quantities on documented household proxies. Emissions are reported at the appliance; outdoor escape and exposure are not estimated. The note supplies a consistent source account for GHG and air applications without converting emission rankings into exposure priorities.
